\hypertarget{chapter-2-memory-types-and-pointers}{%
\section{CHAPTER 2: MEMORY TYPES AND
POINTERS}\label{chapter-2-memory-types-and-pointers}}

\hypertarget{the-address-space-the-map-of-mem-city}{%
\subsubsection{1.1 The Address Space: The Map of
Mem-City}\label{the-address-space-the-map-of-mem-city}}

Every computer program lives in a virtual ``Address Space.'' Imagine
this as a giant, infinite row of mailboxes, each with a unique number
(the address).

\hypertarget{the-layout-of-your-program-in-memory}{%
\paragraph{The Layout of your Program in
Memory}\label{the-layout-of-your-program-in-memory}}

When your program starts, the Operating System divides Mem-City into
several ``Zoning Districts.'' Each district has its own rules, speed
limits, and rent costs.

\begin{longtable}[]{@{}
  >{\raggedright\arraybackslash}p{(\columnwidth - 4\tabcolsep) * \real{0.33}}
  >{\raggedright\arraybackslash}p{(\columnwidth - 4\tabcolsep) * \real{0.33}}
  >{\raggedright\arraybackslash}p{(\columnwidth - 4\tabcolsep) * \real{0.33}}@{}}
\toprule
District & The Zoning Rules & Who Lives Here? \\
\midrule
\endhead
\textbf{The Text Segment} & \textbf{Read-Only Park}: No one is allowed
to change the ground. & Your compiled code (the machine instructions).
It's fixed forever. \\
\textbf{The Data Segment} & \textbf{The Town Square}: Fixed-size statues
that stay forever. & Global variables
(\passthrough{\lstinline!int g = 10;!}) and
\passthrough{\lstinline!static!} variables. \\
\textbf{The BSS Segment} & \textbf{The Empty Lot}: Reserved space for
future statues. & Uninitialized global variables. The OS
zero-initializes these lot for you. \\
\textbf{The Stack} & \textbf{The Quick-Start Desk}: A desk that grows
and shrinks. & Local variables, function parameters, and the ``Return
Address'' (how the CPU knows where to go back to after a function). \\
\textbf{The Heap} & \textbf{The Industrial Warehouse}: Massive space you
rent by the square foot. & Dynamic memory
(\passthrough{\lstinline!new!}, \passthrough{\lstinline!malloc!}). It's
big, but you have to manage it. \\
\bottomrule
\end{longtable}

\begin{center}\rule{0.5\linewidth}{0.5pt}\end{center}

\hypertarget{fireside-chat-why-pointers-break-your-brain}{%
\subsubsection{Fireside Chat: Why Pointers Break Your
Brain}\label{fireside-chat-why-pointers-break-your-brain}}

\textbf{Student}: ``I just don't get it. If I have a variable
\passthrough{\lstinline!int x = 10;!}, why can't I just use
\passthrough{\lstinline!x!}? Why do I need
\passthrough{\lstinline!int* p = \&x;!}?''

\textbf{The Architect}: ``Think about a huge library. If you want to
tell your friend about a great book, you have two choices. 1. You can
photocopy every single page of the book and hand them the pile of paper
(\textbf{Pass by Value}). 2. You can just hand them a slip of paper with
the shelf location: `Floor 2, Row 10, Shelf 4' (\textbf{Pass by
Pointer/Reference}).''

\textbf{Student}: ``Okay, the location is easier. But what if the
librarian moves the book?''

\textbf{The Architect}: ``That's exactly why Pointers are dangerous! If
the book moves but you still have the old address, you're looking at an
empty shelfor worse, a different book entirely. That's a
\textbf{Dangling Pointer}.''

\begin{center}\rule{0.5\linewidth}{0.5pt}\end{center}

\hypertarget{step-by-step-the-life-of-a-pointer}{%
\subsubsection{Step-by-Step: The Life of a
Pointer}\label{step-by-step-the-life-of-a-pointer}}

Let's trace a pointer's life in the CPU registers and RAM.

\begin{lstlisting}[language={C++}]
int main() {
    int secret_number = 42;    // 1. Build a house
    int* spy = &secret_number; // 2. Write down the address
    *spy = 100;                // 3. Go to the address and change the contents
}
\end{lstlisting}

\begin{enumerate}
\def\labelenumi{\arabic{enumi}.}
\tightlist
\item
  \textbf{Step 1}: The CPU asks the OS for 4 bytes on the
  \textbf{Stack}. The OS gives it address
  \passthrough{\lstinline!0x1000!}. The CPU writes the bits for
  \passthrough{\lstinline!42!} into that location.
\item
  \textbf{Step 2}: The CPU asks for another 8 bytes (on a 64-bit system)
  for the pointer \passthrough{\lstinline!spy!}. It stores the value
  \passthrough{\lstinline!0x1000!} into this new house.
\item
  \textbf{Step 3}: The CPU looks at the value in
  \passthrough{\lstinline!spy!} (\passthrough{\lstinline!0x1000!}),
  jumps to that location in RAM, and overwrites the
  \passthrough{\lstinline!42!} with \passthrough{\lstinline!100!}.
\end{enumerate}

\begin{center}\rule{0.5\linewidth}{0.5pt}\end{center}

\hypertarget{common-pointer-street-gangs-traps}{%
\subsubsection{1.2 Common Pointer ``Street Gangs''
(Traps)}\label{common-pointer-street-gangs-traps}}

\begin{longtable}[]{@{}
  >{\raggedright\arraybackslash}p{(\columnwidth - 4\tabcolsep) * \real{0.33}}
  >{\raggedright\arraybackslash}p{(\columnwidth - 4\tabcolsep) * \real{0.33}}
  >{\raggedright\arraybackslash}p{(\columnwidth - 4\tabcolsep) * \real{0.33}}@{}}
\toprule
The Trap & What it is & How to avoid it \\
\midrule
\endhead
\textbf{The Ghost (Wild Pointer)} & A pointer that was never
initialized. It's pointing at a random house in the city. & Always
initialize to \passthrough{\lstinline!nullptr!}. \\
\textbf{The Zombie (Dangling Pointer)} & You deleted the house, but you
still have the address. & Set to \passthrough{\lstinline!nullptr!}
immediately after \passthrough{\lstinline!delete!}. \\
\textbf{The Squatter (Memory Leak)} & You rented a warehouse locker,
threw away the key, and never returned it. & Use \textbf{Smart Pointers}
(RAII). \\
\bottomrule
\end{longtable}

\begin{center}\rule{0.5\linewidth}{0.5pt}\end{center}

\hypertarget{deep-dive-pointer-arithmetic-walking-the-streets}{%
\subsubsection{Deep Dive: Pointer Arithmetic (Walking the
Streets)}\label{deep-dive-pointer-arithmetic-walking-the-streets}}

Pointers are just numbers (addresses), so you can add or subtract from
them. But C++ is smartit knows the ``size'' of the houses.

\begin{itemize}
\tightlist
\item
  If you have an \passthrough{\lstinline!int* p!} pointing at address
  \passthrough{\lstinline!100!}, and you do
  \passthrough{\lstinline!p++!}, it doesn't go to
  \passthrough{\lstinline!101!}.
\item
  It jumps to \passthrough{\lstinline!104!} (because an
  \passthrough{\lstinline!int!} is 4 bytes).
\end{itemize}

\textbf{It's like walking down a street where every house is exactly 4
meters wide. Taking one step forward always puts you at the front door
of the next neighbor.}

\begin{center}\rule{0.5\linewidth}{0.5pt}\end{center}

\hypertarget{memory-types-and-pointers}{%
\section{MEMORY, TYPES, AND POINTERS}\label{memory-types-and-pointers}}

Welcome to the heart of C++. Most languages (Java, Python, JS) try to
hide memory from you. C++ hands you the keys to the city and says,
``Don't burn it down.''

\hypertarget{the-city-of-memory-analogy}{%
\subsubsection{The City of Memory
Analogy}\label{the-city-of-memory-analogy}}

Imagine your computer's RAM is a giant city called \textbf{Mem-City}.

\begin{enumerate}
\def\labelenumi{\arabic{enumi}.}
\tightlist
\item
  \textbf{Memory Addresses}: Every house in Mem-City has a unique street
  address (e.g., \passthrough{\lstinline!0x7ffee6b5a!}).
\item
  \textbf{Variables}: A variable is just a \textbf{House}. When you say
  \passthrough{\lstinline!int x = 5;!}, the Mayor (the OS) builds a
  house, puts the number \passthrough{\lstinline!5!} inside it, and
  names the house ``x''.
\item
  \textbf{Pointers}: A pointer is a \textbf{GPS Device}. It doesn't hold
  a value like \passthrough{\lstinline!5!}; it holds the \textbf{Street
  Address} of a house.
\end{enumerate}

\hypertarget{why-do-we-care}{%
\paragraph{Why do we care?}\label{why-do-we-care}}

In other languages, if you want to give someone your house, you have to
\emph{clone} the entire house and give them the copy. In C++, you just
give them the \textbf{Street Address} (a pointer). Its faster, more
efficient, and allows two people to look at the same house at the same
time.

\begin{center}\rule{0.5\linewidth}{0.5pt}\end{center}

\hypertarget{the-two-districts-stack-vs.-heap}{%
\subsubsection{The Two Districts: Stack
vs.~Heap}\label{the-two-districts-stack-vs.-heap}}

Mem-City is divided into two main districts where variables can live:

\begin{longtable}[]{@{}
  >{\raggedright\arraybackslash}p{(\columnwidth - 6\tabcolsep) * \real{0.25}}
  >{\raggedright\arraybackslash}p{(\columnwidth - 6\tabcolsep) * \real{0.25}}
  >{\raggedright\arraybackslash}p{(\columnwidth - 6\tabcolsep) * \real{0.25}}
  >{\raggedright\arraybackslash}p{(\columnwidth - 6\tabcolsep) * \real{0.25}}@{}}
\toprule
District & Analogy: The Work Space & Lifetime & Speed \\
\midrule
\endhead
\textbf{The Stack} & \textbf{The Desk}: Think of this as your immediate
office desk. You put things on it as you need them. When you leave the
office (function ends), the cleaning crew automatically wipes the desk
clean. & Automatic (ends with \passthrough{\lstinline!\}!}) &
\textbf{Ultra Fast}. Just like grabbing a pen from your desk. \\
\textbf{The Heap} & \textbf{The Warehouse}: A giant storage facility
across town. If you need to store something huge or keep it forever, you
call the Warehouse Manager (\passthrough{\lstinline!new!}) and ask for a
locker. & Manual (You must \passthrough{\lstinline!delete!} it) &
\textbf{Slower}. You have to travel to the warehouse and talk to the
manager. \\
\bottomrule
\end{longtable}

\begin{quote}
\textbf{There are no dumb questions\ldots{}}

\textbf{Q: What happens if I forget to clean out my Warehouse locker
(Heap memory)?} \textbf{A:} You get a \textbf{Memory Leak}. The locker
stays ``rented'' forever, even if your program isn't using it. If you
keep doing this, Mem-City runs out of space and the whole computer
crashes.

\textbf{Q: Why don't I just put everything on the Stack (The Desk)?}
\textbf{A:} Because your desk is small! If you try to put a 1,000-page
book on a tiny desk, you get a \textbf{Stack Overflow}. Use the
Warehouse for the big stuff.
\end{quote}

\begin{center}\rule{0.5\linewidth}{0.5pt}\end{center}

\hypertarget{advanced-pointers-memory}{%
\section{ADVANCED POINTERS \& MEMORY}\label{advanced-pointers-memory}}

\hypertarget{pointer-to-const-vs-const-pointer}{%
\subsection{1.1 Pointer to Const vs Const
Pointer}\label{pointer-to-const-vs-const-pointer}}

\begin{lstlisting}[language={C++}]
#include <iostream>
using namespace std;

int main() {
    int x = 5, y = 10;

    // Pointer to const - can't modify data
    const int* ptr1 = &x;
    // *ptr1 = 10;  // ERROR
    ptr1 = &y;      // OK - can change pointer

    // Const pointer - can't modify pointer
    int* const ptr2 = &x;
    *ptr2 = 10;     // OK - can change data
    // ptr2 = &y;   // ERROR

    // Const pointer to const - can't modify either
    const int* const ptr3 = &x;
    // *ptr3 = 10;  // ERROR
    // ptr3 = &y;   // ERROR

    return 0;
}
\end{lstlisting}

\hypertarget{void-pointers}{%
\subsection{1.2 Void Pointers}\label{void-pointers}}

\begin{lstlisting}[language={C++}]
#include <iostream>
using namespace std;

int main() {
    int x = 42;
    double y = 3.14;

    // Void pointer can point to any type
    void* ptr = &x;
    cout << *(int*)ptr << endl;  // 42

    ptr = &y;
    cout << *(double*)ptr << endl;  // 3.14

    // Generic function using void*
    void print_value(void* ptr, char type) {
        if (type == 'i') {
            cout << *(int*)ptr << endl;
        } else if (type == 'd') {
            cout << *(double*)ptr << endl;
        }
    }

    print_value(&x, 'i');  // 42
    print_value(&y, 'd');  // 3.14

    return 0;
}
\end{lstlisting}

\begin{longtable}[]{@{}
  >{\raggedright\arraybackslash}p{(\columnwidth - 0\tabcolsep) * \real{0.06}}@{}}
\toprule
\endhead
\#\#\# Professional Notes: Pointers \& References \\
\#\#\#\# 1. Arrays as Pointers In C++, an array name decays to a pointer
to its first element in most contexts. * \textbf{Accessing}:
\passthrough{\lstinline!arr[i]!} is equivalent to
\passthrough{\lstinline!*(arr + i)!}. * \textbf{Size}:
\passthrough{\lstinline!sizeof(arr)!} returns the total size in bytes,
whereas \passthrough{\lstinline!sizeof(ptr)!} returns the size of the
pointer (usually 4 or 8 bytes). \\
\#\#\#\# 2. References vs.~Pointers * \textbf{References}: Must be
initialized upon declaration. Cannot be NULL. Cannot be reseated
(re-pointed). * \textbf{Pointers}: Can be initialized later. Can be
NULL. Can point to different objects over time. \\
\textbf{Godhood Tip}: Use references for function parameters to avoid
copying and for operator overloading. Use pointers for dynamic memory
management and optional parameters. \\
\#\#\#\# 3. Pointers to Members Pointers to members are a specialized
feature allowing you to point to a data member or function inside a
class without an instance. ```cpp struct Point \{ int x, y; \}; int
Point::*p\_x = \&Point::x; // Pointer to member x \\
Point p = \{10, 20\}; std::cout \textless\textless{} p.*p\_x
\textless\textless{} std::endl; // Accessing via pointer to member
``` \\
\#\#\#\# 4. The \passthrough{\lstinline!this!} Pointer Inside every
non-static member function, \passthrough{\lstinline!this!} is a hidden
pointer to the current instance. * \textbf{Type}:
\passthrough{\lstinline!T* const!} (or
\passthrough{\lstinline!const T* const!} in const methods). *
\textbf{Usage}: Returning \passthrough{\lstinline!*this!} allows for
method chaining (e.g., in a Fluent API). \\
\bottomrule
\end{longtable}

\hypertarget{professional-notes-language-boundary-storage}{%
\subsubsection{Professional Notes: Language Boundary \&
Storage}\label{professional-notes-language-boundary-storage}}

\hypertarget{c-incompatibilities-the-parents-shadow}{%
\paragraph{1. C Incompatibilities: The Parent's
Shadow}\label{c-incompatibilities-the-parents-shadow}}

While C++ evolved from C, they are distinct languages. *
\textbf{\passthrough{\lstinline!void*!} Conversion}: C allows
\passthrough{\lstinline!int* p = malloc(10);!}. C++ requires
\passthrough{\lstinline!int* p = static\_cast<int*>(malloc(10));!}. *
\textbf{Enumerations}: In C, enums are effectively integers. In C++,
they are distinct types. * \textbf{Tentative Definitions}: C allows
\passthrough{\lstinline!int x; int x;!} at file scope. C++ considers the
second one a redefinition (ODR violation).

\hypertarget{storage-class-specifiers}{%
\paragraph{2. Storage Class Specifiers}\label{storage-class-specifiers}}

\begin{itemize}
\tightlist
\item
  \textbf{\passthrough{\lstinline!static!}}: Internal linkage for
  globals; persistent lifetime for locals.
\item
  \textbf{\passthrough{\lstinline!extern!}}: External linkage. Tells the
  compiler the variable is defined elsewhere.
\item
  \textbf{\passthrough{\lstinline!thread\_local!} (C++11)}: Unique
  instance per thread.
\item
  \textbf{\passthrough{\lstinline!register!}}: (Deprecated in C++11,
  removed in C++17) Hint to use a CPU register.
\item
  \textbf{\passthrough{\lstinline!auto!}}: (C++98) Automatic storage.
  (C++11) Type deduction.
\end{itemize}

\hypertarget{digit-separators-and-binary-literals-c14}{%
\paragraph{3. Digit Separators and Binary Literals
(C++14)}\label{digit-separators-and-binary-literals-c14}}

Use \passthrough{\lstinline!'!} to separate digits for readability:
\passthrough{\lstinline!int x = 1'000'000;!}. Use
\passthrough{\lstinline!0b!} for binary:
\passthrough{\lstinline!int b = 0b1101;!}.

\begin{center}\rule{0.5\linewidth}{0.5pt}\end{center}

\hypertarget{null-pointer-safety}{%
\subsection{1.3 Null Pointer Safety}\label{null-pointer-safety}}

\begin{lstlisting}[language={C++}]
#include <iostream>
using namespace std;

int main() {
    int* ptr = NULL;  // Set to null

    // Always check before dereferencing
    if (ptr != NULL) {
        cout << *ptr << endl;
    } else {
        cout << "Pointer is NULL" << endl;
    }

    // Safer approach
    ptr = new int(42);
    if (ptr) {
        cout << *ptr << endl;
        delete ptr;
        ptr = NULL;
    }

    return 0;
}
\end{lstlisting}

\hypertarget{memory-layout-alignment}{%
\subsection{1.4 Memory Layout \&
Alignment}\label{memory-layout-alignment}}

\begin{lstlisting}[language={C++}]
#include <iostream>
using namespace std;

int main() {
    struct Data {
        char a;     // 1 byte
        int b;      // 4 bytes
        double c;   // 8 bytes
    };

    cout << "Size of Data: " << sizeof(Data) << endl;
    // Likely 16 or 24 (due to alignment padding)

    cout << "Size of char: " << sizeof(char) << endl;      // 1
    cout << "Size of int: " << sizeof(int) << endl;        // 4
    cout << "Size of double: " << sizeof(double) << endl;  // 8

    Data data;
    cout << "Address of a: " << (void*)&data.a << endl;
    cout << "Address of b: " << (void*)&data.b << endl;
    cout << "Address of c: " << (void*)&data.c << endl;

    return 0;
}
\end{lstlisting}

\begin{center}\rule{0.5\linewidth}{0.5pt}\end{center}

\hypertarget{advanced-arrays}{%
\section{ADVANCED ARRAYS}\label{advanced-arrays}}

\hypertarget{dynamic-arrays}{%
\subsection{4.1 Dynamic Arrays}\label{dynamic-arrays}}

\begin{lstlisting}[language={C++}]
#include <iostream>
using namespace std;

int main() {
    // 1D dynamic array
    int size = 5;
    int* arr = new int[size];

    for (int i = 0; i < size; i++) {
        arr[i] = i * 10;
    }

    for (int i = 0; i < size; i++) {
        cout << arr[i] << " ";
    }
    cout << endl;

    delete[] arr;
    arr = NULL;

    // 2D dynamic array
    int rows = 3, cols = 4;
    int** matrix = new int*[rows];
    for (int i = 0; i < rows; i++) {
        matrix[i] = new int[cols];
    }

    // Fill matrix
    for (int i = 0; i < rows; i++) {
        for (int j = 0; j < cols; j++) {
            matrix[i][j] = i * cols + j;
        }
    }

    // Delete matrix
    for (int i = 0; i < rows; i++) {
        delete[] matrix[i];
    }
    delete[] matrix;

    return 0;
}
\end{lstlisting}

\hypertarget{variable-length-arrays-non-standard}{%
\subsection{4.2 Variable Length Arrays
(Non-standard)}\label{variable-length-arrays-non-standard}}

\begin{lstlisting}[language={C++}]
#include <iostream>
using namespace std;

int main() {
    int size;
    cout << "Enter size: ";
    cin >> size;

    // VLA - not standard but supported by many compilers
    int arr[size];  // GCC extension

    for (int i = 0; i < size; i++) {
        arr[i] = i;
    }

    for (int i = 0; i < size; i++) {
        cout << arr[i] << " ";
    }
    cout << endl;

    return 0;
}
\end{lstlisting}

\hypertarget{array-bounds-safety}{%
\subsection{4.3 Array Bounds \& Safety}\label{array-bounds-safety}}

\begin{lstlisting}[language={C++}]
#include <iostream>
using namespace std;

int main() {
    int arr[5] = {10, 20, 30, 40, 50};

    // No bounds checking in C++
    cout << arr[0] << endl;   // 10 (OK)
    cout << arr[10] << endl;  // Undefined behavior!

    // Manual bounds checking
    int index = 5;
    if (index >= 0 && index < 5) {
        cout << arr[index] << endl;
    } else {
        cout << "Index out of bounds" << endl;
    }

    return 0;
}
\end{lstlisting}

\begin{center}\rule{0.5\linewidth}{0.5pt}\end{center}

\hypertarget{const-volatile}{%
\section{CONST \& VOLATILE}\label{const-volatile}}

\hypertarget{const-correctness}{%
\subsection{11.1 Const Correctness}\label{const-correctness}}

\begin{lstlisting}[language={C++}]
#include <iostream>
using namespace std;

int main() {
    int x = 10;

    // const variable
    const int constant = 5;
    // constant = 10;  // ERROR

    // pointer to const
    const int* ptr1 = &x;
    // *ptr1 = 20;  // ERROR
    ptr1 = &constant;  // OK

    // const pointer
    int* const ptr2 = &x;
    *ptr2 = 20;  // OK
    // ptr2 = &constant;  // ERROR

    // const reference
    const int& ref = x;
    // ref = 20;  // ERROR

    cout << x << endl;
    cout << *ptr1 << endl;
    cout << *ptr2 << endl;
    cout << ref << endl;

    return 0;
}
\end{lstlisting}

\hypertarget{volatile-keyword}{%
\subsection{11.2 Volatile Keyword}\label{volatile-keyword}}

\begin{lstlisting}[language={C++}]
#include <iostream>
using namespace std;

int main() {
    // volatile - tells compiler value may change unexpectedly
    volatile int sensor_reading = 0;  // From hardware

    // Compiler won't optimize away reads
    while (sensor_reading < 100) {
        // Check actual value each time, not cached
    }

    // Common use: hardware registers
    volatile int* hardware_register = (volatile int*)0x1000;

    // Each access reads from actual location
    int val1 = *hardware_register;
    int val2 = *hardware_register;

    // Without volatile, compiler might optimize one read

    return 0;
}
\end{lstlisting}

\begin{center}\rule{0.5\linewidth}{0.5pt}\end{center}

\hypertarget{type-casting}{%
\section{TYPE CASTING}\label{type-casting}}

\hypertarget{c-style-casting}{%
\subsection{8.1 C-Style Casting}\label{c-style-casting}}

\begin{lstlisting}[language={C++}]
#include <iostream>
#include <cmath>
using namespace std;

int main() {
    double d = 3.14;

    // C-style cast (avoid in modern C++)
    int i = (int)d;  // 3
    cout << i << endl;

    int x = 65;
    char c = (char)x;  // 'A'
    cout << c << endl;

    float f = (float)d;
    cout << f << endl;

    return 0;
}
\end{lstlisting}

\hypertarget{implicit-conversions}{%
\subsection{8.2 Implicit Conversions}\label{implicit-conversions}}

\begin{lstlisting}[language={C++}]
#include <iostream>
using namespace std;

int main() {
    // Implicit conversions
    int x = 5;
    double d = x;  // int to double (automatic)
    cout << d << endl;  // 5.0

    double d2 = 3.9;
    int y = d2;  // double to int (loses precision)
    cout << y << endl;  // 3

    // Char arithmetic
    char c = 'A';
    int code = c;  // char to int (gets ASCII)
    cout << code << endl;  // 65

    return 0;
}
\end{lstlisting}

\begin{center}\rule{0.5\linewidth}{0.5pt}\end{center}

\hypertarget{enumeration-unions}{%
\section{ENUMERATION \& UNIONS}\label{enumeration-unions}}

\hypertarget{enumerations}{%
\subsection{10.1 Enumerations}\label{enumerations}}

\begin{lstlisting}[language={C++}]
#include <iostream>
using namespace std;

// Enum definition
enum Color { RED, GREEN, BLUE };

enum Direction {
    NORTH = 0,
    EAST = 1,
    SOUTH = 2,
    WEST = 3
};

int main() {
    Color c = RED;
    cout << c << endl;  // 0

    Color colors[3] = {RED, GREEN, BLUE};

    // Switching on enum
    switch (c) {
        case RED:
            cout << "Red" << endl;
            break;
        case GREEN:
            cout << "Green" << endl;
            break;
        case BLUE:
            cout << "Blue" << endl;
            break;
    }

    // Iterate through enum values
    for (int dir = NORTH; dir <= WEST; dir++) {
        cout << "Direction: " << dir << endl;
    }

    return 0;
}
\end{lstlisting}

\hypertarget{unions}{%
\subsection{10.2 Unions}\label{unions}}

\begin{lstlisting}[language={C++}]
#include <iostream>
using namespace std;

// Union - all members share same memory
union Data {
    int i;
    float f;
    char c;
};

int main() {
    Data data;

    cout << "Size of Data: " << sizeof(data) << endl;  // 4 (size of largest member)

    data.i = 10;
    cout << "data.i: " << data.i << endl;     // 10
    cout << "data.f: " << data.f << endl;     // Garbage (overwrites data.i)

    data.f = 3.14;
    cout << "data.i: " << data.i << endl;     // Garbage (overwrites by data.f)
    cout << "data.f: " << data.f << endl;     // 3.14

    // Union useful for memory-constrained systems
    union Variant {
        int int_val;
        double double_val;
        char char_val;
    };

    cout << "Size of Variant: " << sizeof(Variant) << endl;

    return 0;
}
\end{lstlisting}

\begin{center}\rule{0.5\linewidth}{0.5pt}\end{center}

\hypertarget{bitwise-operations}{%
\section{BITWISE OPERATIONS}\label{bitwise-operations}}

\hypertarget{bitwise-operators}{%
\subsection{6.1 Bitwise Operators}\label{bitwise-operators}}

\begin{lstlisting}[language={C++}]
#include <iostream>
using namespace std;

int main() {
    unsigned char a = 5;   // 0101
    unsigned char b = 3;   // 0011

    // AND
    cout << (a & b) << endl;  // 0001 = 1

    // OR
    cout << (a | b) << endl;  // 0111 = 7

    // XOR
    cout << (a ^ b) << endl;  // 0110 = 6

    // NOT (bitwise complement)
    cout << (~a) << endl;     // 1010 = 250 (for unsigned char)

    // Left shift
    cout << (a << 1) << endl; // 1010 = 10

    // Right shift
    cout << (b >> 1) << endl; // 0001 = 1

    return 0;
}
\end{lstlisting}

\hypertarget{bit-manipulation-techniques}{%
\subsection{6.2 Bit Manipulation
Techniques}\label{bit-manipulation-techniques}}

\begin{lstlisting}[language={C++}]
#include <iostream>
using namespace std;

int main() {
    unsigned int num = 5;  // 0101

    // Check if bit is set
    int bit_pos = 2;
    bool is_set = (num >> bit_pos) & 1;
    cout << "Bit " << bit_pos << " is: " << is_set << endl;

    // Set a bit
    num |= (1 << 1);  // Set bit 1
    cout << "After setting bit 1: " << num << endl;  // 7 (0111)

    // Clear a bit
    num &= ~(1 << 1);  // Clear bit 1
    cout << "After clearing bit 1: " << num << endl;  // 5 (0101)

    // Toggle a bit
    num ^= (1 << 0);  // Toggle bit 0
    cout << "After toggling bit 0: " << num << endl;  // 4 (0100)

    // Count set bits
    unsigned int count = 0;
    unsigned int temp = num;
    while (temp) {
        count += temp & 1;
        temp >>= 1;
    }
    cout << "Number of set bits: " << count << endl;

    return 0;
}
\end{lstlisting}

\begin{center}\rule{0.5\linewidth}{0.5pt}\end{center}

\hypertarget{professional-notes-chapter-8-arrays}{%
\section{Professional Notes: Chapter 8:
Arrays}\label{professional-notes-chapter-8-arrays}}

Arrays are elements of the same type placed in adjoining memory
locations. The elements can be individually referenced by a unique
identier with an added index. This allows you to declare multiple
variable values of a specic type and access them individually without
needing to declare a variable for each value. Section 8.1: Array
initialization An array is just a block of sequential memory locations
for a specic type of variable. Arrays are allocated the same way as
normal variables, but with square brackets appended to its name {[}{]}
that contain the number of elements that t into the array memory. The
following example of an array uses the typ int, the variable name
arrayOfInts, and the number of elements {[}5{]} that the array has space
for: int arrayOfInts{[}5{]}; An array can be declared and initialized at
the same time like this int arrayOfInts{[}5{]} = \{10, 20, 30, 40, 50\};
When initializing an array by listing all of its members, it is not
necessary to include the number of elements inside the square brackets.
It will be automatically calculated by the compiler. In the following
example, it's 5: int arrayOfInts{[}{]} = \{10, 20, 30, 40, 50\}; It is
also possible to initialize only the rst elements while allocating more
space. In this case, dening the length in brackets is mandatory. The
following will allocate an array of length 5 with partial
initialization, the compiler initializes all remaining elements with the
standard value of the element type, in this case zero. int
arrayOfInts{[}5{]} = \{10,20\}; // means 10, 20, 0, 0, 0 Arrays of other
basic data types may be initialized in the same way. char
arrayOfChars{[}5{]}; // declare the array and allocate the memory, don't
initialize char arrayOfChars{[}5{]} = \{ `a', `b', `c', `d', `e' \} ;
//declare and initialize double arrayOfDoubles{[}5{]} = \{1.14159,
2.14159, 3.14159, 4.14159, 5.14159\}; string arrayOfStrings{[}5{]} = \{
``C++'', ``is'', ``super'', ``duper'', ``great!''\}; It is also
important to take note that when accessing array elements, the array's
element index(or position) starts from 0. int array{[}5{]} = \{
10/\emph{Element no.0}/, 20/\emph{Element no.1}/, 30, 40,
50/\emph{Element no.4}/\}; std::cout \textless\textless{} array{[}4{]};
//outputs 50 std::cout \textless\textless{} array{[}0{]}; //outputs 10
Section 8.2: A xed size raw array matrix (that is, a 2D raw array) // A
fixed size raw array matrix (that is, a 2D raw array). \#include
\#include using namespace std; auto main() -\textgreater{} int \{ int
const n\_rows = 3; int const n\_cols = 7; int const
m{[}n\_rows{]}{[}n\_cols{]} = // A raw array matrix. \{ \{ 1, 2, 3, 4,
5, 6, 7 \}, \{ 8, 9, 10, 11, 12, 13, 14 \}, \{ 15, 16, 17, 18, 19, 20,
21 \} \}; for( int y = 0; y \textless{} n\_rows; ++y ) \{ for( int x =
0; x \textless{} n\_cols; ++x ) \{ cout \textless\textless{} setw( 4 )
\textless\textless{} m{[}y{]}{[}x{]}; // Note: do NOT use m{[}y,x{]}! \}
cout \textless\textless{} `\n'; \} \} Output: 1 2 3 4 5 6 7 8 9 10 11 12
13 14 15 16 17 18 19 20 21 C++ doesn't support special syntax for
indexing a multi-dimensional array. Instead such an array is viewed as
an array of arrays (possibly of arrays, and so on), and the ordinary
single index notation {[}i{]} is used for each level. In the example
above m{[}y{]} refers to row y of m, where y is a zero-based index. Then
this row can be indexed in turn, e.g.~m{[}y{]}{[}x{]}, which refers to
the xth item or column of row y. I.e. the last index varies fastest, and
in the declaration the range of this index, which here is the number of
columns per row, is the last and innermost size specied. Since C++
doesn't provide built-in support for dynamic size arrays, other than
dynamic allocation, a dynamic size matrix is often implemented as a
class. Then the raw array matrix indexing notation m{[}y{]}{[}x{]} has
some cost, either by exposing the implementation (so that e.g.~a view of
a transposed matrix becomes practically impossible) or by adding some
overhead and slight inconvenience when it's done by returning a proxy
object from operator{[}{]}. And so the indexing notation for such an
abstraction can and will usually be dierent, both in look-and-feel and
in the order of indices, e.g.~m(x,y) or m.at(x,y) or m.item(x,y).
Section 8.3: Dynamically sized raw array // Example of raw dynamic size
array. It's generally better to use std::vector. \#include // std::sort
\#include using namespace std; auto int\_from( istream\& in )
-\textgreater{} int \{ int x; in \textgreater\textgreater{} x; return x;
\} auto main() -\textgreater{} int \{ cout \textless\textless{}
``Sorting n integers provided by you.\textbackslash n''; cout
\textless\textless{} ``n?''; int const n = int\_from( cin ); int* a =
new int{[}n{]}; // Allocation of array of n items. for( int i = 1; i
\textless= n; ++i ) \{ cout \textless\textless{} ``The \#''
\textless\textless{} i \textless\textless{} " number, please: ``;
a{[}i-1{]} = int\_from( cin ); \} sort( a, a + n ); for( int i = 0; i
\textless{} n; ++i ) \{ cout \textless\textless{} a{[}i{]}
\textless\textless{} ' `; \} cout
\textless\textless{}'\textbackslash n'; delete{[}{]} a; \} A program
that declares an array T a{[}n{]}; where n is determined a run-time, can
compile with certain compilers that support C99 variadic length arrays
(VLAs) as a language extension. But VLAs are not supported by standard
C++. This example shows how to manually allocate a dynamic size array
via a new{[}{]}-expression, int* a = new int{[}n{]}; // Allocation of
array of n items. then use it, and nally deallocate it via a
delete{[}{]}-expression: delete{[}{]} a; The array allocated here has
indeterminate values, but it can be zero-initialized by just adding an
empty parenthesis (), like this: new int\href{}{n}. More generally, for
arbitrary item type, this performs a value-initialization. As part of a
function down in a call hierarchy this code would not be exception safe,
since an exception before the delete{[}{]} expression (and after the
new{[}{]}) would cause a memory leak. One way to address that issue is
to automate the cleanup via e.g.~a std::unique\_ptr smart pointer. But a
generally better way to address it is to just use a std::vector: that's
what std::vector is there for. Section 8.4: Array size: type safe at
compile time \#include // size\_t, ptrdiff\_t
//----------------------------------- Machinery: using Size =
ptrdiff\_t; template\textless{} class Item, size\_t n \textgreater{}
constexpr auto n\_items( Item (\&){[}n{]} ) noexcept -\textgreater{}
Size \{ return n; \} //----------------------------------- Usage:
\#include using namespace std; auto main() -\textgreater{} int \{ int
const a{[}{]} = \{3, 1, 4, 1, 5, 9, 2, 6, 5, 4\}; Size const n =
n\_items( a ); int b{[}n{]} = \{\}; // An array of the same size as a.
(void) b; cout \textless\} The C idiom for array size,
sizeof(a)/sizeof(a{[}0{]}), will accept a pointer as argument and will
then generally yield an incorrect result. For C++11 using C++11 you can
do: std::extent\textless decltype(MyArray)\textgreater::value; Example:
char MyArray{[}{]} = \{ `X',`o',`c',`e' \}; const auto n =
std::extent\textless decltype(MyArray)\textgreater::value; std::cout
\textless\textless{} n \textless\textless{}''\n``; // Prints 4 Up till
C++17 (forthcoming as of this writing) C++ had no built-in core language
or standard library utility to obtain the size of an array, but this can
be implemented by passing the array by reference to a function template,
as shown above. Fine but important point: the template size parameter is
a size\_t, somewhat inconsistent with the signed Size function result
type, in order to accommodate the g++ compiler which sometimes insists
on size\_t for template matching. With C++17 and later one may instead
use std::size, which is specialized for arrays. Section 8.5: Expanding
dynamic size array by using std::vector // Example of std::vector as an
expanding dynamic size array. \#include // std::sort \#include \#include
// std::vector using namespace std; int int\_from( std::istream\& in )
\{ int x = 0; in \textgreater\textgreater{} x; return x; \} int main()
\{ cout \textless\textless{}''Sorting integers provided by you.\n``;
cout \textless\textless{}''You can indicate EOF via F6 in Windows or
Ctrl+D in Unix-land.\n``; vector a; // Zero size by default. while( cin
) \{ cout \textless\textless{}''One number, please, or indicate EOF: ";
int const x = int\_from( cin ); if( !cin.fail() ) \{ a.push\_back( x );
\} // Expands as necessary. \} sort( a.begin(), a.end() );  int const n
= a.size(); for( int i = 0; i \textless{} n; ++i ) \{ cout
\textless\textless{} a{[}i{]} \textless\textless{} ' `; \} cout
\textless\textless{}'\n'; \} std::vector is a standard library class
template that provides the notion of a variable size array. It takes
care of all the memory management, and the buer is contiguous so a
pointer to the buer (e.g.~\&v{[}0{]} or v.data()) can be passed to API
functions requiring a raw array. A vector can even be expanded at run
time, via e.g.~the push\_back member function that appends an item. The
complexity of the sequence of n push\_back operations, including the
copying or moving involved in the vector expansions, is amortized O(n).
Amortized: on average. Internally this is usually achieved by the vector
doubling its buer size, its capacity, when a larger buer is needed. E.g.
for a buer starting out as size 1, and being repeatedly doubled as
needed for n=17 push\_back calls, this involves 1 + 2 + 4 + 8 + 16 = 31
copy operations, which is less than 2n = 34. And more generally the sum
of this sequence can't exceed 2n. Compared to the dynamic size raw array
example, this vector-based code does not require the user to supply (and
know) the number of items up front. Instead the vector is just expanded
as necessary, for each new item value specied by the user. Section 8.6:
A dynamic size matrix using std::vector for storage Unfortunately as of
C++14 there's no dynamic size matrix class in the C++ standard library.
Matrix classes that support dynamic size are however available from a
number of 3rd party libraries, including the Boost Matrix library (a
sub-library within the Boost library). If you don't want a dependency on
Boost or some other library, then one poor man's dynamic size matrix in
C++ is just like vector\textless vector\textgreater{} m( 3, vector( 7 )
); where vector is std::vector. The matrix is here created by copying a
row vector n times where n is the number of rows, here 3. It has the
advantage of providing the same m{[}y{]}{[}x{]} indexing notation as for
a xed size raw array matrix, but it's a bit inecient because it involves
a dynamic allocation for each row, and it's a bit unsafe because it's
possible to inadvertently resize a row. A more safe and ecient approach
is to use a single vector as storage for the matrix, and map the client
code's (x, y) to a corresponding index in that vector: // A dynamic size
matrix using std::vector for storage.
//--------------------------------------------- Machinery: \#include //
std::copy \#include // assert \#include // std::initializer\_list
\#include // std::vector \#include // ptrdiff\_t namespace my \{ using
Size = ptrdiff\_t; using std::initializer\_list; using std::vector;
template\textless{} class Item \textgreater{} class Matrix \{ private:
vector items\_; Size n\_cols\_; auto index\_for( Size const x, Size
const y ) const -\textgreater{} Size \{ return y\emph{n\_cols\_ + x; \}
public: auto n\_rows() const -\textgreater{} Size \{ return
items\_.size()/n\_cols\_; \} auto n\_cols() const -\textgreater{} Size
\{ return n\_cols\_; \} auto item( Size const x, Size const y )
-\textgreater{} Item\& \{ return items\_{[}index\_for(x, y){]}; \} auto
item( Size const x, Size const y ) const -\textgreater{} Item const\& \{
return items\_{[}index\_for(x, y){]}; \} Matrix(): n\_cols\_( 0 ) \{\}
Matrix( Size const n\_cols, Size const n\_rows ) : items\_(
n\_cols}n\_rows ) , n\_cols\_( n\_cols ) \{\} Matrix(
initializer\_list\textless{} initializer\_list \textgreater{} const\&
values ) : items\_() , n\_cols\_( values.size() == 0? 0 :
values.begin()-\textgreater size() ) \{ for( auto const\& row : values )
\{ assert( Size( row.size() ) == n\_cols\_ ); items\_.insert(
items\_.end(), row.begin(), row.end() ); \} \} \}; \} // namespace my
//--------------------------------------------- Usage: using my::Matrix;
auto some\_matrix() -\textgreater{} Matrix \{ return \{ \{ 1, 2, 3, 4,
5, 6, 7 \}, \{ 8, 9, 10, 11, 12, 13, 14 \}, \{ 15, 16, 17, 18, 19, 20,
21 \} \}; \} \#include \#include using namespace std; auto main()
-\textgreater{} int \{ Matrix const m = some\_matrix(); assert(
m.n\_cols() == 7 ); assert( m.n\_rows() == 3 ); for( int y = 0, y\_end
= m.n\_rows(); y \textless{} y\_end; ++y ) \{ for( int x = 0, x\_end =
m.n\_cols(); x \textless{} x\_end; ++x ) \{ cout \textless{} Note: not
\passthrough{\lstinline!m[y][x]!}! \} cout \textless\} \} Output: 1 2 3
4 5 6 7 8 9 10 11 12 13 14 15 16 17 18 19 20 21 The above code is not
industrial grade: it's designed to show the basic principles, and serve
the needs of students learning C++. For example, one may dene operator()
overloads to simplify the indexing notation.
